\begin{abstract}
We study whether oracle foreground masks improve feature-based fine-grained fungal recognition when strong frozen visual features are already available. On a 215-class closed-set subset of FungiTastic, we extract frozen DINOv3 ViT-7B descriptors and compare seven token constructions at `224' and `448' resolution: CLS, mean pooled patches, mask-aware pooled patches, and simple concatenations of global and local features. All variants are evaluated with the same small MLP classifier to isolate the effect of the representation. We find that masks help most when used as a complementary local cue rather than as a standalone descriptor. Fusing the CLS token with mask-aware pooled patches gives the best accuracy at both resolutions, reaching \textbf{74.04}\% at `224' and \textbf{75.53}\% at `448', improving over plain CLS by 1.82 and 1.23 points, respectively. In contrast, naive mean pooling over all patches is consistently weak, and branch-wise input normalization does not help. These results suggest that, for frozen fungal features, oracle masks are most useful for refining local evidence while preserving global context.
\end{abstract}
